% Statistics source: 2007湖南文卷原始 PDF 第1页第7题；矢量路径中每 8.5041 pt 对应 0.5% 频率组距。
% solid segments: contiguous one-metre histogram bin outlines [30,31),...,[50,51).
% dashed guides: frequency-density reference lines 0.5%, 1%, 2%; labels avoid bars and axes.
% Relations: each rectangle height is the frequency density for its water-level interval;
% the original chart's tail-frequency question uses these contiguous bins.
\begin{tikzpicture}
  \begin{axis}[exam statistical histogram,
    width=16.0cm,
    height=11.665cm,
    xmin=29.6,
    xmax=51.8,
    ymin=0,
    ymax=8.4,
    axis lines=left,
    axis line style={-Stealth},
    tick style={black},
    ylabel={},
    xtick={30,31,32,33,48,49,50,51},
    ytick={0},
    yticklabels={0},
    xticklabel style={font=\normalsize,anchor=north},
    yticklabel style={font=\normalsize},
    extra y ticks={0.5,1,2},
    extra y tick labels={0.5\%,1\%,2\%},
    extra y tick style={
      grid=major,
      grid style={dashed,black!70},
      tick style={draw=none},
      yticklabel style={font=\normalsize},
    },
    enlarge x limits=false,
    clip=false,
  ]
    \examstatxlabel{水位（米）};
    \addplot[
      ybar interval,
      draw=black,
      fill=none,
      line width=0.4pt,
    ] coordinates {
      (30,2.5) (31,3.5) (32,4) (33,5) (34,5.5) (35,6.5)
      (36,5.5) (37,7.5) (38,7) (39,6.5) (40,6) (41,5)
      (42,5.5) (43,5) (44,6) (45,5) (46,5.5) (47,4.5)
      (48,2) (49,1) (50,0.5) (51,0)
    };
    \examstatylabel{\(\dfrac{\text{频率}}{\text{组距}}\)};
  \end{axis}
\end{tikzpicture}
